\chapter{Run Khi Lên Bảng}
\chapterintrobi{Ở dưới còn biết, lên bảng lại quên sạch}{Shake when standing at the board}{Không ít bạn có kiến thức trong đầu nhưng đánh rơi gần hết trên quãng đường từ bàn tới bảng.}

\begin{lucbatbox}{Lục bát: Từ bàn tới bảng mất chữ}
Ở dưới em nhớ rất nhanh\
Lên trên bảng viết bỗng thành mong manh\
Phấn kia cầm cũng không đành\
Tay run, đầu trắng như cành trước mưa\
Bạn nhìn còn thấy thương chưa\
Chỉ vì hồi hộp mà vừa quên ngang
\end{lucbatbox}

\begin{tutuyetbox}{Thất ngôn tứ tuyệt: Bước lên là mất sóng}
Ngồi dưới em còn khá tự tin\
Lên bảng xong rồi chữ bỗng im\
Nếu phần bình tĩnh theo lên trước
Bài kia đâu hóa cuộc đi tìm
\end{tutuyetbox}

\begin{ghichu}
Run khi lên bảng không có nghĩa là không biết. Rất nhiều khi chỉ là áp lực đứng trước ánh nhìn làm trí nhớ bị khóa tạm thời.
\end{ghichu}

\begin{englishnote}
\textbf{at the board}: trên bảng.\par
\textbf{freeze}: đứng hình.\par
\textbf{I knew it at my seat, but I froze at the board.}: Em biết khi ngồi dưới nhưng lên bảng lại đứng hình.
\end{englishnote}